\subsection{The 1842 Apportionment Act}

The Apportionment Act of June~25, 1842 \citep{apportionment1842}
is notable for two reasons: it was the first to mandate
single-member congressional districts, and it used Webster's
major fractions method.
The 1840 Census had produced a controversy: under Jefferson's method
(still in use), several large states received seat counts that seemed
disproportionately high.
Webster, who was then Secretary of State, advocated for his method
as a fairer alternative.

The 1842 Act apportioned 223 seats to 26 states.
The total House size was reduced from 242 (the 1830 apportionment)
to 223, reflecting concerns about the growing size of the House.
Webster's method was used for this apportionment, though the Act
was controversial and some states initially refused to hold
single-member-district elections.

\subsection{The 1850 and Hamilton Interlude}

The Apportionment Act of 1850 replaced Webster with Hamilton's
largest-remainder method for apportionments through 1900.
The political motivation was that Hamilton's quota-satisfying
property was seen as more equitable: every state would receive either
its lower or upper quota, with no state receiving more or fewer seats
than its ``exact'' entitlement.
Webster was not used again until 1911.

\subsection{The 1911 Act and Webster's Return}

The Apportionment Act of 1911 \citep{apportionment1911} returned to
Webster's method, partly because Hamilton's Alabama Paradox
had become politically embarrassing.
The 1900 Census had demonstrated the Population Paradox under Hamilton,
where Virginia lost a seat to Maine despite Virginia's population
growing faster.
Webster, as a divisor method, was paradox-immune and was seen as a
partial solution.

The 1910 apportionment allocated 433 seats (the House size was increased
from 391 to 433, then to 435 in 1913 to account for the new states
of Arizona and New Mexico).
Webster produced seat counts that were widely regarded as fair
and that did not trigger paradox complaints.

\subsection{The 1930 Non-Apportionment and the 1941 Transition}

The 1930 Census was the only census since 1790 after which no
apportionment occurred.
Congress could not agree on the apportionment method (Webster vs.\
a version of Hamilton) or the House size.
The political deadlock partly reflected the fact that Webster's method,
applied to the 1930 Census, would have shifted seats from rural to
urban states, which was politically unacceptable to the rural-dominated
Congress of the era.

The Reapportionment Act of 1929 \citep{apportionment1929} mandated
automatic apportionment after each census but did not specify the method.
This ambiguity led to the 1941 Act, which explicitly adopted
Huntington-Hill and ended the legislative debate over apportionment method.
Webster was not seriously reconsidered for federal use after 1941,
though it remains the dominant method in comparative electoral systems.
